% =====================================================================
% Note 7, size buckets. Moved out of the Supplementary Methods, 24 November (guide's
% comment: no Supplementary Methods; NMI allows algorithm detail as a Note).
% =====================================================================

\section{Size bucketing of gene groups}
\label{note:buckets}

Each gene group is read by its own transformer, with its own weights. At a
token width of 32 every matrix product inside one of these transformers is
small enough that launching it costs more than computing it, so running the
44 groups one after another leaves the processor mostly idle. We therefore
stacked the groups' weights and computed several groups in one batched
call. A batched call needs one sequence length, so every group in it is
padded with empty positions up to the longest. Attention costs grow with
the square of sequence length, so the padding is expensive: run as a single
block, 46.7\% of the attention arithmetic would be spent on empty positions
(Supplementary Table~\ref{stab:buckets}).

We instead sorted the groups by size and divided them into four buckets,
each padded only to its own longest group. The division is the one that
minimises the total padded attention work,
\begin{equation}
W = \sum_{b} n_b L_b^{2},
\label{eq:bucketwork}
\end{equation}
where $n_b$ is the number of groups in bucket $b$ and $L_b$ its padded
length, the size of its largest group plus the one learned summary token.
Once the groups are sorted by size, an optimal division can always be taken
as contiguous runs of the sorted list, so the exact minimum of
equation~(\ref{eq:bucketwork}) is found by a dynamic program over the cut
points rather than approximated by splitting the list evenly. In the
reported vocabulary the four buckets hold 9, 14, 13 and 8 groups, and the
share of arithmetic spent on padding falls to 12.9\%. Each further bucket
would remove less than the one before.

Bucketing changes the order in which groups are computed and nothing else.
Each group keeps its own weights. Padded positions are masked as keys, so
no gene attends to them, and they are never read as outputs, so the eight
pooled summary statistics are computed over a group's real genes only. A
test in the code base copies one set of weights into the bucketed and the
unbucketed encoder and requires the two to give the same output for every
group. The implementation can also drop the mask, and with it reach the
fused attention kernel \citep{dao2022flashattention}, for a bucket whose groups are all one size. No
bucket in the reported vocabulary is uniform, so every bucket carries a
mask and the saving is in arithmetic alone.

\begin{table}[t]
\centering
\caption[Size buckets]{\textbf{The size buckets of the reported model.}
Padded length includes the learned summary token. Arithmetic on padding is
the share of padded attention work, $n_b L_b^2$ summed, spent on empty
positions. Generated by running the model's own partition on its
vocabulary, and checked against the partition the training run recorded.}
\label{stab:buckets}
\small
% GENERATED by figures/scripts/supp_methods.py from bucketing.plan_buckets
% and vocabulary_v1.json; asserted against the training log. Do not edit.
\begin{tabular}{@{}lrrrr@{}}
\toprule
Bucket & Groups & Genes per group & Padded length & Arithmetic on padding \\
\midrule
\quad 1 & 9 & 247--265 & 266 & 6.0\% \\
\quad 2 & 14 & 283--343 & 344 & 15.7\% \\
\quad 3 & 13 & 355--422 & 423 & 14.5\% \\
\quad 4 & 8 & 436--493 & 494 & 10.9\% \\
\addlinespace
\quad All 4 buckets & 44 & 247--493 & & 12.9\% \\
\quad One block, for comparison & 44 & 247--493 & 494 & 46.7\% \\
\midrule
\multicolumn{5}{@{}l}{\textit{Other bucket counts}} \\
\quad 2 buckets & & & & 26.8\% \\
\quad 3 buckets & & & & 18.0\% \\
\quad 5 buckets & & & & 9.6\% \\
\quad 6 buckets & & & & 7.5\% \\
\quad 8 buckets & & & & 4.6\% \\
\bottomrule
\end{tabular}

\end{table}
